\chapter{結論}\label{chap:conclusion}

% てきとうに散らかしてあるだけ
本稿では，従来の価値反復ROSパッケージの走り出すまでにかかる
時間をA*を併用することで短縮する手法を提案した．
実装するA*について，$xy$平面を探索するA*と，
$xy\theta$空間を探索するA*の2種類のA*
を実装し，その効果をシミュレータを用いて検証した．
その結果，A*の計算結果が十分に小さかった場合，
走り出しまでの時間が短縮されることが分かった．
その一方で，A*の計算に時間がかかる場合には，
状態価値関数に分岐が生成され，ロボットの
移動に時間がかかってしまう現状を確認した．

A*探索の計算時間の増加と移動時間の短縮の効果のバランスから，
$xy\theta$を探索するA*より，$xy$平面を探索するA*
のほうが走り出しまでの時間を短縮するための
探索手法としては適していた．

今後は，より環境の大域的な構造を捉え，
複数の解像度でA*探索を行う方法や，
作成した一般ボロノイ図をトポロジカルマップとして
探索を行う方法を検証する．

